% Background
\section{Background}
\label{sec:background}

\subsection{Eternal inflation and pocket-domain counting}

In slow-roll inflation, quantum fluctuations of the inflaton field can cause
different spatial regions to undergo different numbers of $e$-folds, producing
a fractal-like structure of pocket domains \cite{guth_eternal}.
Because the overall volume grows without bound, any counting over pocket
domains requires a regularization---a measure---to obtain finite
probabilities \cite{freivogel_measure}.

\subsection{The measure problem}

A global census that simply counts trajectories (or pocket domains) equally
suffers from severe trajectory-cutoff dependence. As the number of sampled
trajectories increases, the distribution shifts toward whichever basin has the
largest population. Different cutoff choices produce different answers, and
there is no natural stopping point.

This problem is well known in the literature \cite{freivogel_measure} and
motivates proposals such as the causal diamond measure, the scale-factor
cutoff, and the light-cone cutoff \cite{winitzki_measure,garriga_vilenkin,bousso_causal}.

\subsection{Coarse-graining analogy}

The situation is structurally reminiscent of renormalization in quantum field
theory: a cutoff is introduced to regularize divergences, and physical
predictions must be insensitive to the cutoff scale.
In the eternal-inflation context, a good measure should be stable under
increases in the trajectory cutoff $N$.
\ISR draws on the coarse-graining intuition that unobserved degrees
of freedom can be averaged over, producing effective weights for
the observed sector.
This is a structural analogy, not a Wilsonian RG derivation: the
scale $\ell$ is a locality scale in sector space, not a momentum scale,
and no decoupling theorem is assumed.

\subsection{Configuration-space locality prior}\label{sec:locality_prior}

Locality here does not mean causal signaling between bubbles.
It means proximity in the shared stochastic landscape used to generate
sectors. The prior is justified only inside a model where sectors arise
from common inflaton dynamics.

ISR assumes that, conditional on a shared inflationary landscape, sectors
near the observable patch in configuration space are more relevant to the
local effective ensemble than sectors arbitrarily far away.
This is a modeling prior, not an observable causal relation.

Basin boundaries are exactly where this prior is expected to fail.
ISR is valuable only if it exposes those failures, not if it hides them.

Throughout this work, ``locality'' denotes proximity in the configuration space
of a specified stochastic landscape. It does not denote causal contact, signal
propagation.

\subsection{Relation to stochastic inflation and open-system thinking}

Stochastic inflation \cite{starobinsky_stochastic,burgess_stochastic} treats the inflaton as
a classical field driven by quantum noise.
Our simulation uses a Langevin-like evolution in a multi-well potential,
similar in spirit to open-system treatments where the environment
(short-wavelength modes) affects the long-wavelength field.
The unobserved-sector perspective extends this by treating entire unobserved
inflationary domains as an effective environment.

\subsection{Relation to existing measure proposals}

Table~\ref{tab:measure_comparison} situates ISR among existing cosmological
measure proposals. ISR is not a replacement for spacetime-volume measures;
it is a conditional weighting scheme inside a generated sector ensemble.

\begin{table*}[t]
  \centering
  \caption{Comparison of ISR with existing cosmological measure proposals. ISR does not regulate spacetime volume divergences; it is a conditional weighting scheme on a generated sector ensemble.}
  \label{tab:measure_comparison}
  \small
  \begin{tabular}{lcccc}
    \toprule
    Measure & Regulates volume? & Regulates trajectory cutoff? & Tested in & Relation to ISR \\
    \midrule
    Causal diamond \cite{winitzki_measure} & Yes & No & No & Composable: ISR weights could be applied inside \\
    Scale-factor cutoff \cite{winitzki_measure} & Yes & No & No & Composable: ISR weights could be applied inside \\
    Light-cone cutoff \cite{garriga_vilenkin} & Yes & No & No & Composable: ISR weights could be applied inside \\
    Global census & No & No (counts equally) & This work & Baseline: ISR is compared against it \\
    ISR (this work) & No & Yes & This work & --- \\
    \bottomrule
  \end{tabular}
\end{table*}

